\subsection{Concluding Remarks}
\label{sec:conclusion}

\noindent
EMT is an OS framework for developing and evaluating
    OS memory management for new memory translation architectures.
Our work shows the importance of understanding
    hardware translation schemes from the OS perspective.
Specifically, we show that fast translation schemes like ECPT 
    can incur new challenges for the OS.
%    which can be overlooked with hardware-centric design process.
% EMT aims to enable such experience
%    for systems with hardware-software co-design.
We will release and make EMT and the emulator toolchains 
    as an open platform and encourage new hardware architectures
    to experiment with modern OSes.
% Our ultimate goal is an infrastructure 
%   where new translation architectures
%    can be fairly compared and comprehensively understood. 
% The aspiration of EMT is beyond an experiment platform. 
With the significant diversity of emerging workloads
    and increasing heterogeneity of interconnected memory devices,
    it becomes harder to foresee a one-size-fits-all translation scheme.
%    as shown by recent research on hybrid schemes.
Hence, OS extensibility for different translation schemes 
    is critical to enable % user-space,
    specialized translation.
% Memory translation has been a basic building block of 
%   heterogeneous computing devices such as IOMMU~\cite{Amit2010},
%    NeuMMU~\cite{hyun2020neummu}, and MPU~\cite{picorel2017near};
%    some are also known to suffer from 
%    off-TLB translation overhead.
% EMT currently focuses on memory translation for CPU, but
We hope that EMT design starts a practical 
    journey towards extensible OS kernels for translation of heterogeneous devices.

%\newpage

%Based on traditional IOMMU, NPU MMU incorporates pending request merging buffer
%    to merge identical page translations,
%    more page table walkers to improve MMU throughput,
%    and page table walk registers (lightweight cache) to reduce the number of page table walks.

% need to integrate into the middle
\begin{comment}
Due to the greater complexity of code and logic, macro expansion techniques are
difficult to apply in high-level functions. Furthermore, we observe that
high-level functions are rarely replaced on a large scale in the vast majority
of memory translation designs. This is because these high-level functions are
inherently more architecture-independent. So, it may not be worth it to refactor
every function call like what we do for the low layers. For this reason, we
employ low-overhead runtime checks as shown in Figure \ref{code:if} to implement
function override for high-level interface implementations. Specifically, we
define a set of VFS-like function pointer structures in units of \igrp, allowing
the user to assign alternative implementations for each function. Then, in the
Linux function of the same name, we insert a lightweight check at the function
head that will use the alternative implementation to process the request when
present. Otherwise, the default Linux implementation will work as-is.

This approach proves to be low overhead due to the presence of hardware branch predictors.
Since function pointers do not change often, such predictors can eliminate the checking overhead with high efficiency.
Hence, we accomplish on-demand functions redirection with near-zero overhead,
protecting performance while satisfying the \interface design.
\end{comment}